\statement{B.76}{Refuted}{Retained-only special collisions do not guarantee relative deletion-kernel surjectivity}
The following proposed shortcut to the source condition in B.75 is
false, both discretely and for the actual images in full profinite
completions: if a special fiber has only collisions among retained
punctures, while the punctures to be filled remain distinct and finite,
then its deletion kernel surjects onto the ambient deletion kernel.
The failure persists with constant fiber degree and a continuous
section over the entire base.

\begin{proof}
\emph{The actual surfaces and based comparison.}\par
Let \(B=\mathbb C\setminus\{1\}\), and put
\[
\begin{aligned}
 Y&=\{(\lambda,z)\in B\times\mathbb C:
                       z\notin\{0,\lambda,1\}\},\\
 X&=\{(\lambda,z)\in B\times\mathbb C:
                       z\notin\{0,\lambda\}\}.
\end{aligned}
 \tag{B.76.1}
\]
The inclusion \(j:Y\hookrightarrow X\) fills exactly
\(U=\{z=1\}\cap X\). The special fibers over \(\lambda=0\)
are \(L_Y=\mathbb C\setminus\{0,1\}\) and
\(L_X=\mathbb C\setminus\{0\}\). We use \(z_*=2\) and
\((\lambda_*,z_*)=(0,2)\) throughout. The four inclusion maps
form a strictly commuting based square. Its two deletion kernels are
\[
\begin{aligned}
 K_L&=\ker\bigl(\pi_1(L_Y,2)\longrightarrow\pi_1(L_X,2)\bigr),\\
 K_Y&=\ker\bigl(\pi_1(Y,(0,2))\longrightarrow
                         \pi_1(X,(0,2))\bigr).
\end{aligned}
 \tag{B.76.2}
\]
All groups in this display are discrete topological fundamental
groups. We will exhibit an element of \(K_Y\) outside even the
profinite closure of the full special-fiber image.

At \(\lambda=0\), only the retained punctures \(0,\lambda\)
collide. The filled puncture \(1\) remains distinct. The polynomial
\(z(z-\lambda)(z-1)\) has constant degree three. The continuous
section \(z=2+|\lambda|\) avoids every deleted point on all of
\(B\), for both surfaces. Thus all the proposed shortcut's local
conditions hold.

\emph{The total group and the special source.}\par
There is an algebraic product chart with its explicit inverse:
\[
 (\lambda,z)\longmapsto
       \left(z,r=\frac{\lambda-z}{1-z}\right),
 \qquad \lambda=z+r(1-z).
 \tag{B.76.3}
\]
The conditions on \(Y\) become exactly \(z,r\notin\{0,1\}\).
Consequently
\[
       Y\simeq(\mathbb C\setminus\{0,1\})^2,
       \qquad\pi_1(Y,(0,2))\simeq F\times F,
       \quad F=F_2.
 \tag{B.76.4}
\]
The basepoint becomes \((2,2)\). In this chart the special
fiber is the graph of
\[
                         f(z)=\frac{z}{z-1}.
 \tag{B.76.5}
\]
This is an involutory M\"obius automorphism of the three-punctured
sphere fixing \(2\). Hence the actual special-fiber group image is
the graph of the based isomorphism \(f_*:F\to F\).
Identifying the second factor by \(f_*^{-1}\), this image is
\(H=\{(g,g):g\in F\}\). In particular the image of \(K_L\)
is contained in this precise subgroup \(H\), without taking its
normal closure.

\emph{The filled group.}\par
We calculate \(\pi_1(X)\) directly. Project to
\(\lambda\in\mathbb C\setminus\{1\}\). Away from zero the
fiber has meridians \(p,q\) about \(0,\lambda\). A loop about
\(\lambda=0\) gives the two-puncture full twist. Filling this
base disk kills its section lift and imposes precisely
\([p,q]=1\): the full twist acts by conjugation by \(pq\),
and triviality of that action is equivalent to this commutator
relation. This is the ordinary node form of the based van Kampen
calculation. A loop about the excluded base value \(1\) has
trivial braid action, since the moving puncture's corresponding
loop contracts in \(\mathbb C\setminus\{0\}\). The far section
adds no basepoint winding. There are no further finite
degenerations or degree drops. Thus
\[
 \pi_1(X)=\langle p,q,t\mid
       [p,q]=[t,p]=[t,q]=1\rangle\simeq\mathbb Z^3.
 \tag{B.76.6}
\]
This also proves the particular three-line arrangement calculation
being used: the deleted lines \(\lambda=1,z=0,z=\lambda\)
have three distinct ordinary nodes.

Choose a free basis \(a,b\) of the first factor in (B.76.4).
Since (B.76.6) is abelian, the nonidentity element
\[
                              w=([a,b],1)
 \tag{B.76.7}
\]
belongs to \(K_Y\). It is outside the graph \(H\), whose
intersection with the first factor is trivial. Hence the map
\(K_L\to K_Y\) is not onto.

\emph{A finite certificate and the full completed conclusion.}\par
Send the identified free basis in each factor to
\(a\mapsto(12),b\mapsto(23)\). This defines an actual
epimorphism
\[
              \rho:\pi_1(Y)\twoheadrightarrow S_3\times S_3.
 \tag{B.76.8}
\]
The graph \(H\) maps to diagonal \(S_3\), whereas \(w\) maps
to a nonidentity three-cycle in the first factor and identity
in the second. The exact chart identities, the order \(36\),
the diagonal order \(6\), and this nonmembership are verified in
\path{code/check_special_fiber_relative_counterexample.py}; its
genuine output and full certificate are
\path{records/check_special_fiber_relative_counterexample.out} and
\path{records/special_fiber_relative_counterexample.json}.
The six cosets of diagonal \(S_3\) give a connected finite-cover
certificate containing the source image and missing \(w\).

The continuous extension of (B.76.8) separates \(w\) from the
closure of \(H\) in \(\widehat{\pi_1(Y)}\). Its image in
\(\widehat{\pi_1(X)}\) is nevertheless one, because its
discrete image already is one. Therefore it is also outside the
image of the completed source deletion kernel, while belonging
to the ambient completed deletion kernel. This argument uses
the actual maps of completions and a finite separating quotient;
it does not assume that an infinitely generated kernel has its
full induced topology.

The omitted base value \(\lambda=1\) carries a triple involving
\(U\) and retained divisors. Its monodromy survives as a base
action; it is not a relation imposed by filling a base disk.
This explains why the retained-only collision at the selected
fiber cannot justify the proposed conclusion. In the actual
icosahedral ratio pencil, the corresponding excluded-base actions
and their common based transport must be retained. The present
counterexample does not determine that actual source image and
does not supply a nonidentity element of \(\mathcal K\).
\end{proof}
